\documentclass[conference]{IEEEtran}

\usepackage[T1]{fontenc}
\usepackage{times}
\usepackage{graphicx}
\usepackage{amsmath,amssymb}
\usepackage{booktabs}
\usepackage{url}
\usepackage[hidelinks]{hyperref}

\pagestyle{plain}

\title{An API Security and Anomaly Detection Gateway for\\
Controlled Request Evaluation}

\author{\IEEEauthorblockN{Student One, Student Two, and Student Three}
\IEEEauthorblockA{Department of Computer Science\\
University Name\\
City, State, USA\\
\{student.one, student.two, student.three\}@university.edu}}

\begin{document}
\maketitle
\thispagestyle{plain}

\begin{abstract}
This paper presents a placeholder description of an API security and anomaly
detection gateway. The gateway evaluates requests before they reach a protected
application, combines deterministic rules with an extensible anomaly-detector
interface, and records sanitized audit events. The report is intentionally
formatted as an IEEE-style student submission and contains placeholder prose
that can be replaced with project-specific results. The proposed architecture
separates transport, normalization, detection, decision policy, and persistence
so that each concern can be tested independently. Evaluation is described using
latency percentiles, throughput, and confusion-matrix measures.
\end{abstract}

\begin{IEEEkeywords} % see rules for alphabetical order, quantity of words, etc.
API security, anomaly detection, request validation, audit logging, gateway,
security evaluation
\end{IEEEkeywords}

\section{Introduction}
Application programming interfaces are a primary boundary between clients and
modern services. A request that is syntactically valid may still violate an
endpoint policy, contain suspicious input, or exhibit behavior inconsistent
with normal traffic. This project investigates a defensive gateway that creates
an explicit decision before forwarding traffic to a protected API.

The motivating idea is simple: a gateway should make security decisions using
a bounded and normalized representation of a request. It should also explain
the decision sufficiently for operators to investigate it, while avoiding the
storage of credentials and complete payloads. This concern is consistent with
the principle of minimizing sensitive data in operational systems
\cite{nist80053,owaspapi}.

The remainder of this paper is a report template with placeholder text. The
architecture is presented in Section II, the request-processing design in
Section III, the evaluation approach in Section IV, and limitations and future
work in Sections V and VI.

\section{System Architecture}
The gateway is organized as a modular monolith. A single deployable service
contains interfaces between ingress, normalization, detector implementations,
decision policy, and audit persistence. This arrangement reduces deployment
complexity while preserving boundaries that can later support separate
components. Figure~\ref{fig:pipeline} summarizes the intended data flow.

\begin{figure}[!t]
\centering
\fbox{\begin{minipage}{0.94\columnwidth}
\centering
\small
Client $\rightarrow$ Ingress $\rightarrow$ Validation $\rightarrow$
Normalization $\rightarrow$ Detectors $\rightarrow$ Policy\\[2pt]
\hfill $\searrow$ allow $\rightarrow$ Protected API\\
\hfill $\searrow$ block $\rightarrow$ Safe error\\
\hfill $\searrow$ audit event $\rightarrow$ PostgreSQL
\end{minipage}}
\caption{Conceptual request-processing pipeline.}
\label{fig:pipeline}
\end{figure}

The ingress layer assigns a request identifier, applies transport limits, and
returns consistent status codes. Authentication and schema validation occur
before forwarding. The normalizer converts framework-specific data into an
internal record containing the method, route template, identity metadata,
selected header attributes, query metadata, and a bounded body representation.
Secrets are removed before detector or audit code receives the record.

The detector pipeline provides two complementary paths. Deterministic rules
recognize known malformed inputs and policy violations. An adapter interface
allows a statistical or machine-learning detector to return a score,
threshold, model version, and classification without coupling the transport
layer to a particular framework. The policy layer combines these results into
an explicit \textit{allow}, \textit{block}, or future \textit{review} decision.

\section{Request Processing and Data Protection}
The processing sequence is designed to make the security boundary observable.
First, the gateway receives a request and assigns a monotonic start time. Next,
it authenticates the caller and validates endpoint-specific fields. Requests
that fail validation are rejected before reaching a protected service. The
normalizer then creates a bounded record for the detection pipeline. This
record is intentionally smaller than the original request.

An audit event stores the request identifier, timestamp, method, route
template, decision, downstream status, detector classification, reason code,
bounded latency, and an optional irreversible correlation identifier. It does
not store authorization headers, cookies, passwords, API keys, access tokens,
complete request bodies, or raw uploaded files. These choices reduce the
consequences of an audit-store compromise and support privacy-aware operations
\cite{nistprivacy}.

For a request $r$, let $D(r)$ be the detector result and $P(D(r))$ the policy
decision. A simple threshold model can be represented as
\begin{equation}
  P(D(r)) =
  \begin{cases}
    \text{block}, & s(r) \geq \tau,\\
    \text{allow}, & s(r) < \tau,
  \end{cases}
  \label{eq:threshold}
\end{equation}
where $s(r)$ is a detector score and $\tau$ is a configured threshold. A
production implementation may add rule precedence, endpoint-specific policy,
or an analyst review path. The equation is included as a compact description
of the decision concept rather than as a claim about final project results.

\begin{table}[!t]
\caption{Illustrative Gateway Measurements}
\label{tab:metrics}
\centering
\begin{tabular}{@{}lll@{}}
\toprule
Metric & Purpose & Example target \\
\midrule
P50 latency & Typical overhead & $<10$ ms \\
P95 latency & Tail overhead & $<50$ ms \\
Throughput & Capacity & $>500$ req/s \\
False-positive rate & User impact & $<2\%$ \\
\bottomrule
\end{tabular}
\end{table}

\section{Evaluation Methodology}
Evaluation should use versioned synthetic data and authorized test traffic.
The benign portion should represent ordinary endpoint usage, while the test
portion should contain controlled examples of malformed requests and policy
violations. Each run should record the detector configuration, threshold,
model version, hardware, dataset composition, and known limitations. This
documentation makes comparisons more meaningful and avoids presenting a single
number without context.

The primary performance measures are gateway-added latency at the 50th, 95th,
and 99th percentiles, throughput, CPU utilization, memory utilization, and
audit-write outcomes. Detection quality is summarized with true positives,
false positives, true negatives, and false negatives. Precision and recall are
defined as
\begin{equation}
  \text{precision}=\frac{TP}{TP+FP}, \qquad
  \text{recall}=\frac{TP}{TP+FN}.
  \label{eq:quality}
\end{equation}

The following placeholder text should be replaced by the group's measured
results. Lorem ipsum dolor sit amet, consectetur adipiscing elit. Integer
vehicula, justo at feugiat posuere, neque turpis volutpat erat, at elementum
libero nibh sed nunc. Curabitur non sem vel arcu tincidunt consequat. Donec
tempor, massa vitae dignissim commodo, justo urna malesuada lectus, at posuere
magna lectus non nibh. These sentences are deliberately generic and should not
be interpreted as empirical evidence.

\section{Implementation Considerations}
The prototype uses a Python web service and a relational database for
structured audit records. Containerized local development makes it possible to
run the gateway and database with repeatable configuration. Runtime settings
such as credentials, thresholds, and model parameters should be supplied
through the deployment environment rather than embedded in source code.

\subsection{Deployment Configuration}
Configuration should be provided at runtime and reviewed separately from the
application source. This subsection demonstrates the expected IEEE heading
hierarchy.

\subsubsection{Local Development}
Local development uses repeatable container configuration and synthetic data.
This subsubsection demonstrates a third-level heading. Replace this placeholder
text with the team's actual setup and experimental procedure.

Failure behavior is part of the security design. Malformed requests fail
validation and are not forwarded. Deterministic rule matches produce a block
decision. Detector timeouts and dependency failures follow an explicit policy
and are observable as operational events. An audit-store interruption must not
silently change a blocked request into an allowed request or cause credentials
to be written to fallback logs.

Lorem ipsum dolor sit amet, consectetur adipiscing elit. Praesent consequat
diam non est gravida, at tincidunt mi faucibus. Suspendisse potenti. Aenean
porta, nibh vitae consequat laoreet, justo massa lacinia eros, in efficitur
libero neque sit amet dolor. Sed vulputate sem at augue commodo, in fermentum
libero cursus. The placeholder paragraph illustrates the expected density of
body text and can be replaced with implementation details, screenshots, or
experiment observations.

\section{Limitations and Future Work}
The gateway is an experimental security control and does not replace secure
application authorization, network controls, or independently validated
production defenses. A detector can miss novel behavior, and an aggressive
threshold can reject legitimate requests. Evaluation data may also fail to
represent traffic encountered after deployment. These limitations should be
reported alongside any performance or accuracy claim.

Future work includes endpoint-aware baselines, drift monitoring, stronger
configuration validation, a controlled analyst-review workflow, and a more
systematic comparison of rule and statistical detectors. The project can also
measure the operational cost of audit persistence and test recovery behavior
under dependency failures. Lorem ipsum dolor sit amet, consectetur adipiscing
elit. Aliquam erat volutpat. Morbi in lorem sit amet elit accumsan consequat.

\section{Conclusion}
This paper described an IEEE-style placeholder report for an API security and
anomaly detection gateway. The central design separates request handling,
normalization, detection, policy, and audit persistence. That separation makes
the system easier to test and allows future detector implementations without
changing the transport contract. The final submission should replace the
placeholder prose with the team's authorship, implementation evidence,
measurements, and discussion of limitations.

\section*{Acknowledgment}
The authors may use this section to acknowledge instructors, teammates, or
resources that supported the project.

\bibliographystyle{IEEEtran}
\bibliography{references}

\end{document}
